\documentclass[a4paper,twoside]{article}

\usepackage{morewrites}

\usepackage[svgnames,dvipsnames,x11names]{xcolor}
\usepackage{graphicx}
\usepackage[english]{babel}
\usepackage[colorlinks=true, linkcolor=NavyBlue, urlcolor=NavyBlue, citecolor=NavyBlue]{hyperref}
\usepackage[a4paper]{geometry}
\usepackage{ts-fonts}
\usepackage{booktabs}
\usepackage{tabularx}
\usepackage{amsmath}
\usepackage{float}
\usepackage{seqsplit}
\newcolumntype{Y}{>{\centering\arraybackslash}X}
\newcolumntype{Z}{>{\raggedleft\arraybackslash}X}
\usepackage{ts-callouts}
\usepackage{ts-code}

\usepackage{lastpage}
\usepackage{refcount}
\makeatletter
\def\ps@plain{
  \let\@oddhead\@empty
  \let\@evenhead\@empty
  \def\@oddfoot{
    \hfil
    \ifnum\getpagerefnumber{LastPage}>1\relax
      \thepage
    \fi
    \hfil}
  \let\@evenfoot\@oddfoot
}
\makeatother

\title{Supported Markdown Syntax}
\author{}\date{}

\begin{document}

\maketitle

\thispagestyle{plain}
TeXSmith bundles Python-Markdown together with a curated set of PyMdown
extensions. The combination lets you author rich documentation while keeping
output predictable for \LaTeX{} conversion. This page summarizes the syntax you can
use out of the box and points to the extension behind each feature.

\begin{callout}[callout info]{Renderer defaults}
The CLI and API both enable the exact same extension list defined in
\texttt{texsmith.adapters.markdown.DEFAULT\_MARKDOWN\_EXTENSIONS}. You can always
override the list with CLI flags or API options, but the features below are
available without additional configuration.
\end{callout}
\section{Core Markdown (Always On)}\label{core-markdown-always-on}

\begin{code}{markdown}{}{}
\begin{Verbatim}[commandchars=\\\{\},breaklines, breakanywhere, commandchars=\\\{\}]
\PY{g+gh}{\PYZsh{} Heading 1}
\PY{g+gu}{\PYZsh{}\PYZsh{} Heading 2}
\PY{g+gu}{\PYZsh{}\PYZsh{}\PYZsh{} Heading 3}

\PY{g+gs}{**bold**}, \PY{g+ge}{*italic*}, \PY{g+gd}{\PYZti{}\PYZti{}strikethrough\PYZti{}\PYZti{}}, \PY{l+s+sb}{`inline code`}

[\PY{n+nt}{Links}](\PY{n+na}{https://www.example.com}) and ![\PY{n+nt}{images}](\PY{n+na}{assets/logo.svg})

\PY{k}{\PYZgt{} }\PY{g+ge}{Blockquote text}

\PY{k}{\PYZhy{}}\PY{+w}{ }Unordered item
\PY{+w}{  }\PY{k}{\PYZhy{}}\PY{+w}{ }Nested item
\PY{k}{1.} Ordered item
\PY{k}{2.} Next item

\PY{g+gu}{Horizontal rules:}
\PY{g+gu}{\PYZhy{}\PYZhy{}\PYZhy{}}
\end{Verbatim}
\end{code}
All standard Markdown constructs---headings, emphasis, lists, code blocks,
blockquotes, links, images, and horizontal rules---render exactly as you would
expect. In \LaTeX{} output, horizontal rules become \texttt{\textbackslash{}clearpage} page breaks rather
than a literal line. TeXSmith relies on fenced code blocks by default, so triple
backticks (\texttt{``}) are the recommended way to author code samples.

\section{Extension Cheat Sheet}\label{extension-cheat-sheet}

\begin{center}
\begin{tabularx}{\linewidth}{>{\raggedright\arraybackslash}X>{\raggedright\arraybackslash}X>{\raggedright\arraybackslash}X>{\raggedright\arraybackslash}X}
\toprule
\textbf{Feature} & \textbf{Extension} & \textbf{Package} & \textbf{Example} \\

\midrule
Definition lists & \texttt{def\_list} & \texttt{markdown} & \texttt{Term: Definition} \\
Footnotes & \texttt{footnotes} & \texttt{markdown} & \texttt{Footnote ref1} \\
Abbreviations & \texttt{abbr} & \texttt{markdown} & \texttt{*[HTML]: HyperText Markup Language} \\
Admonitions & \texttt{admonition} & \texttt{markdown} & \texttt{!!! note  Body} \\
Attribute lists & \texttt{attr\_list} & \texttt{markdown} & \texttt{![Alt](image.png)\{ width="50\%" \}} \\
Tables & \texttt{tables} & \texttt{markdown} & Pipe-delimited tables \\
Markdown in HTML & \texttt{md\_in\_html} & \texttt{markdown} & Markdown inside custom \texttt{<div>} blocks \\
SmartyPants & \texttt{pymdownx.smartsymbols} & \texttt{pymdown-\allowbreak{}extensions} & Auto-converts quotes/dashes \\
Highlighted code & \texttt{pymdownx.highlight} & \texttt{pymdown-\allowbreak{}extensions} & Adds syntax highlighting + anchors \\
Inline highlighting & \texttt{pymdownx.inlinehilite} & \texttt{pymdown-\allowbreak{}extensions} & \texttt{print("hi")} \\
Details/summary & \texttt{pymdownx.details} & \texttt{pymdown-\allowbreak{}extensions} & \texttt{???+ note Title} \\
SuperFences & \texttt{pymdownx.superfences} & \texttt{pymdown-\allowbreak{}extensions} & Nest code fences safely \\
Task lists & \texttt{pymdownx.tasklist} & \texttt{pymdown-\allowbreak{}extensions} & \texttt{-\allowbreak{} [x] Done} \\
Better emphasis & \texttt{pymdownx.betterem} & \texttt{pymdown-\allowbreak{}extensions} & Fixes edge cases with underscores \\
MagicLink & \texttt{pymdownx.magiclink} & \texttt{pymdown-\allowbreak{}extensions} & Autolinks URLs/issues \\
Keys & \texttt{pymdownx.keys} & \texttt{pymdown-\allowbreak{}extensions} & \texttt{Ctrl+Alt+Del} \\
Tabbed content & \texttt{pymdownx.tabbed} & \texttt{pymdown-\allowbreak{}extensions} & Content tabs \\
Snippets & \texttt{pymdownx.snippets} & \texttt{pymdown-\allowbreak{}extensions} & Include external Markdown snippets \\
Caret markup & \texttt{pymdownx.caret} & \texttt{pymdown-\allowbreak{}extensions} & \texttt{insert} \\
Mark (highlight) & \texttt{pymdownx.mark} & \texttt{pymdown-\allowbreak{}extensions} & \texttt{highlight} \\
Tilde syntax & \texttt{pymdownx.tilde} & \texttt{pymdown-\allowbreak{}extensions} & Subscript / superscript \\
Critic markup & \texttt{pymdownx.critic} & \texttt{pymdown-\allowbreak{}extensions} & Editorial annotations \\
Emoji & \texttt{pymdownx.emoji} & \texttt{pymdown-\allowbreak{}extensions} & \texttt{:sparkles:} or \texttt{:fontawesome-\allowbreak{}regular-\allowbreak{}face-\allowbreak{}smile:} \\
Fancy lists & \texttt{pymdownx.fancylists} & \texttt{pymdown-\allowbreak{}extensions} & Extended list markers \\
Blocks caption & \texttt{pymdownx.blocks.caption} & \texttt{pymdown-\allowbreak{}extensions} & Captions for fenced blocks \\
Blocks HTML & \texttt{pymdownx.blocks.html} & \texttt{pymdown-\allowbreak{}extensions} & Named block wrappers \\
Snippets of LaTeX & \texttt{texsmith.extensions.latex\_raw} & bundled & Raw LaTeX fence \\
Missing footnotes guard & \texttt{texsmith.extensions.missing\_footnotes} & bundled & Warns when references lack definitions \\

\bottomrule
\end{tabularx}
\end{center}

Want a generated table of contents? Add the Python-Markdown \texttt{toc} extension with \texttt{-\allowbreak{}x toc} or \texttt{-\allowbreak{}-\allowbreak{}enable-\allowbreak{}extension toc}---it is no longer enabled by default.

Use the table above as a quick pointer. The following sections provide more
context and runnable examples.

\section{Working with Admonitions}\label{working-with-admonitions}

\begin{code}{markdown}{}{}
\begin{Verbatim}[commandchars=\\\{\},breaklines, breakanywhere, commandchars=\\\{\}]
!!! warning \PYZdq{}LaTeX toolchain\PYZdq{}
    Remember to install TeX Live, MiKTeX, or MacTeX before running \PY{l+s+sb}{`texsmith \PYZhy{}\PYZhy{}build`}.
\end{Verbatim}
\end{code}
Admonitions render as highlighted callouts in HTML and as tcolorbox blocks in the
\LaTeX{} output. Combine them with tabs or details blocks to create layered
walkthroughs.

\section{Tables and Definition Lists}\label{tables-and-definition-lists}

\begin{code}{markdown}{}{}
\begin{Verbatim}[commandchars=\\\{\},breaklines, breakanywhere, commandchars=\\\{\}]
| Option | Description |
| \PYZhy{}\PYZhy{}\PYZhy{}\PYZhy{}\PYZhy{}\PYZhy{} | \PYZhy{}\PYZhy{}\PYZhy{}\PYZhy{}\PYZhy{}\PYZhy{}\PYZhy{}\PYZhy{}\PYZhy{}\PYZhy{}\PYZhy{} |
| \PY{l+s+sb}{`\PYZhy{}\PYZhy{}list\PYZhy{}extensions`} | Prints enabled Markdown extensions |
| \PY{l+s+sb}{`\PYZhy{}\PYZhy{}debug`} | Shows full tracebacks |

Term
: Definition content
\end{Verbatim}
\end{code}
Tables use the Python-Markdown \texttt{tables} extension while definition lists come
from \texttt{def\_list}. Both convert cleanly into \LaTeX{} environments.

\section{Task Lists and Checkboxes}\label{task-lists-and-checkboxes}

\begin{code}{markdown}{}{}
\begin{Verbatim}[commandchars=\\\{\},breaklines, breakanywhere, commandchars=\\\{\}]
\PY{k}{\PYZhy{} }\PY{k}{[x]} Validate MkDocs navigation
\PY{k}{\PYZhy{} }\PY{k}{[ ]} Document template slots
\end{Verbatim}
\end{code}
Task lists automatically render checkboxes in HTML. In \LaTeX{} they become custom
itemize entries with inline symbols.

\section{Keyboard Shortcuts}\label{keyboard-shortcuts}

\begin{code}{markdown}{}{}
\begin{Verbatim}[commandchars=\\\{\},breaklines, breakanywhere, commandchars=\\\{\}]
Use ++ctrl+s++ to save changes and ++ctrl+shift+b++ to build the docs.
\end{Verbatim}
\end{code}
The \texttt{pymdownx.keys} extension turns the markup above into keyboard glyphs, which
carry across to PDFs through the TeXSmith formatter.

\section{Embedding Raw LaTeX}\label{embedding-raw-latex}

Use the \texttt{/// latex} fence when native \LaTeX{} is required:

\begin{code}{markdown}{}{}
\begin{Verbatim}[commandchars=\\\{\},breaklines, breakanywhere, commandchars=\\\{\}]
/// latex
\PYZbs{}begin\PYZob{}align\PYZcb{}
E \PYZam{}= mc\PYZca{}2 \PYZbs{}\PYZbs{}
\PYZbs{}nabla \PYZbs{}cdot \PYZbs{}vec\PYZob{}E\PYZcb{} \PYZam{}= \PYZbs{}frac\PYZob{}\PYZbs{}rho\PYZcb{}\PYZob{}\PYZbs{}varepsilon\PYZus{}0\PYZcb{}
\PYZbs{}end\PYZob{}align\PYZcb{}
///
\end{Verbatim}
\end{code}
TeXSmith passes the block straight to the renderer, letting you mix handcrafted
\LaTeX{} with converted Markdown content. For inline adjustments, drop
\texttt{\{latex\}[commands]} right into the paragraph:

\begin{code}{markdown}{}{}
\begin{Verbatim}[commandchars=\\\{\},breaklines, breakanywhere, commandchars=\\\{\}]
The chapter ends here \PYZob{}latex\PYZcb{}[\PYZbs{}clearpage] before appendices.
\end{Verbatim}
\end{code}
\section{Snippet Includes}\label{snippet-includes}

The \texttt{pymdownx.snippets} extension lets you avoid duplication:

\begin{code}{markdown}{}{}
\begin{Verbatim}[commandchars=\\\{\},breaklines, breakanywhere, commandchars=\\\{\}]
\PYZhy{}\PYZhy{}8\PYZlt{}\PYZhy{}\PYZhy{} \PYZdq{}includes/built\PYZhy{}in\PYZhy{}tasks.md\PYZdq{}
\end{Verbatim}
\end{code}
Create an \texttt{includes} directory under \texttt{docs/} and share fragments across pages.
The same mechanism can pull example Markdown from the samples used in automated
tests, keeping docs and fixtures aligned.

\section{When You Need More}\label{when-you-need-more}

\begin{itemize}
\item{} Use \texttt{texsmith -\allowbreak{}-\allowbreak{}list-\allowbreak{}extensions} to see the live extension list.
\item{} Disable or add extensions via the \texttt{-\allowbreak{}-\allowbreak{}enable-\allowbreak{}extension} and \texttt{-\allowbreak{}-\allowbreak{}disable-\allowbreak{}extension}
  flags in \texttt{texsmith} or through \texttt{ConversionRequest.markdown\_extensions}
  in the API.
\item{} If a feature relies on a third-party executable (for example Mermaid to PDF),
  make sure the binary is available on the build worker before running
  \texttt{texsmith -\allowbreak{}-\allowbreak{}build}.

\end{itemize}
With these extensions enabled, TeXSmith can faithfully render everything from
simple README-style guides to complex, reference-heavy manuals.

\end{document}
